\documentclass[11pt,a4paper]{article}

% ── Codificación, fuentes y lenguaje ───────────────────────────────────────────
\usepackage[utf8]{inputenc}
\usepackage[T1]{fontenc}
\usepackage[default,scale=0.95]{sourcesanspro}
\renewcommand{\familydefault}{\sfdefault}
\usepackage[spanish,es-noshorthands,es-tabla]{babel}

% ── Geometría y espaciado ─────────────────────────────────────────────────────
\usepackage[top=2.2cm, bottom=2.5cm, left=2.2cm, right=2.2cm, headheight=15pt]{geometry}
\usepackage{setspace}
\setstretch{1.18}
\usepackage{parskip}

% ── Matemáticas y unidades ────────────────────────────────────────────────────
\usepackage{amsmath}
\usepackage{amssymb}
\usepackage{siunitx}

% ── Circuitos ─────────────────────────────────────────────────────────────────
\usepackage[european, straightvoltages]{circuitikz}

% ── Tablas ────────────────────────────────────────────────────────────────────
\usepackage{booktabs}
\usepackage{tabularx}
\usepackage{array}
\usepackage{longtable}
\usepackage{colortbl}
\usepackage{enumitem}

% ── Colores, cajas e iconos ───────────────────────────────────────────────────
\usepackage[dvipsnames]{xcolor}
\usepackage{tcolorbox}
\tcbuselibrary{skins, breakable}
\usepackage{fontawesome5}
\usepackage{graphicx}

% ── Listados de código (dark-mode infográfico) ────────────────────────────────
%   Estilo base: fondo oscuro con números de línea. Los bloques lstlisting se
%   envuelven automáticamente en una caja tcolorbox redondeada.
\usepackage{listings}

% ── Secciones con barra de color (infografía) ────────────────────────────────
\usepackage{titlesec}
\titleformat{\section}
  {\normalfont\LARGE\bfseries\color{azulNoche}}
  {\colorbox{cyanNeon}{\color{fondoOscuro}\thesection}}{0.7em}{}
  [\vspace{2pt}\color{cyanNeon}\rule{\linewidth}{1.8pt}]
\titlespacing*{\section}{0pt}{24pt}{10pt}
\titleformat{\subsection}
  {\normalfont\large\bfseries\color{azulNoche}}
  {\textcolor{cyanNeon}{\thesubsection}}{0.7em}{}
  [\vspace{1pt}\color{grisLinea}\rule{\linewidth}{0.6pt}]
\titlespacing*{\subsection}{0pt}{14pt}{6pt}
\titleformat{\subsubsection}
  {\normalfont\normalsize\bfseries\color{azulNoche}}
  {\textcolor{cyanNeon}{\thesubsubsection}}{0.7em}{}
\titlespacing*{\subsubsection}{0pt}{10pt}{4pt}

% ── Cabeceras y pies de página ────────────────────────────────────────────────
\usepackage{fancyhdr}
\usepackage{lastpage}
\pagestyle{fancy}
\fancyhf{}
\renewcommand{\headrulewidth}{0pt}
\renewcommand{\footrulewidth}{0pt}
\fancyfoot[C]{%
  \begin{minipage}{\linewidth}
    \vspace{2pt}
    \color{cyanNeon}\rule{\linewidth}{1.2pt}\\[2pt]
    \centering\color{grisTexto}\footnotesize
    Página \thepage\ de \pageref{LastPage}
  \end{minipage}%
}

% ── Hiperenlaces ──────────────────────────────────────────────────────────────
\usepackage[colorlinks=true, linkcolor=azulNoche, urlcolor=cyanNeon]{hyperref}
\sloppy
\emergencystretch 3em

% ── Paleta de colores — tecnológico dark-mode ──────────────────────────────────
\definecolor{fondoOscuro}{HTML}{0D1B2A}%             Dark navy — fondo banda título
\definecolor{verdeTurquesa}{HTML}{004D40}%           Verde turquesa — bandas de marca
\definecolor{azulNoche}{HTML}{1B3A6B}%               Midnight blue — secciones, marcos
\definecolor{azulMedio}{HTML}{2A5298}%               Azul medio — variante de acento
\definecolor{cyanNeon}{HTML}{00C8E0}%                Cyan eléctrico — acentos primarios
\definecolor{verdeSignal}{HTML}{00B887}%             Verde señal — fortalezas, éxito
\definecolor{naranjaVivo}{HTML}{FF7043}%             Naranja vivo — mejoras, advertencias
\definecolor{rojoAlerta}{HTML}{E53935}%              Rojo alerta — errores
\definecolor{grisPapel}{HTML}{F4F6FA}%               Fondo tarjetas claras
\definecolor{grisLinea}{HTML}{C8D0E0}%               Bordes sutiles
\definecolor{grisTexto}{HTML}{6B7A99}%               Texto secundario
\definecolor{amarilloNota}{HTML}{FFB300}%            Notas / advertencias doradas
\definecolor{codigoFondo}{HTML}{1E2D3D}%             Fondo del bloque de código
\definecolor{codigoComentario}{HTML}{7A8FA6}%        Comentarios
\definecolor{codigoCadena}{HTML}{FF9D5C}%            Cadenas / strings
\definecolor{codigoKeyword}{HTML}{00C8E0}%           Palabras clave
\definecolor{codigoNumero}{HTML}{00E5A0}%            Números

% ── Banda de título: portada infográfica ──────────────────────────────────────
%   Diseño en 2 capas: banda de color (por defecto fondo oscuro) + franja de
%   acento cyan abajo. Uso: \bandaTitulo[color]{título}{subtítulo}.
%   El icono se puede cambiar con \renewcommand{\iconoBanda}{...} (FontAwesome).
\newcommand{\iconoBanda}{microchip}
\newcommand{\bandaTitulo}[3][fondoOscuro]{%
  \begin{tcolorbox}[
    enhanced,
    colback=#1, colframe=#1,
    boxrule=0pt, arc=8pt, outer arc=8pt,
    width=\linewidth,
    top=18pt, bottom=6pt, left=14pt, right=14pt,
    borderline south={3.5pt}{0pt}{cyanNeon},
  ]
    \begin{minipage}[c]{0.07\linewidth}
      \centering
      \textcolor{cyanNeon}{\fontsize{30}{30}\selectfont\faIcon{\iconoBanda}}
    \end{minipage}%
    \hspace{8pt}%
    \begin{minipage}[c]{0.88\linewidth}
      {\fontsize{20}{24}\selectfont\bfseries\color{white}#2}\par\vspace{5pt}%
      \textcolor{cyanNeon!80!white}{\small\faIcon{angle-right}~#3}%
    \end{minipage}
    \vspace{4pt}
  \end{tcolorbox}%
  \vspace{8pt}%
}

% ── Tarjetas de datos (infografía) ────────────────────────────────────────────
\tcbset{
  tarjetaDato/.style={
    enhanced, breakable,
    arc=6pt, outer arc=6pt,
    colback=grisPapel, colframe=azulNoche,
    boxrule=1pt,
    fonttitle=\bfseries\small, coltitle=white,
    attach boxed title to top left={yshift=-3mm, xshift=6mm},
    boxed title style={
      colback=azulNoche, colframe=azulNoche,
      arc=4pt, boxrule=0pt,
    },
    drop fuzzy shadow=azulNoche!30!white,
    top=8pt, bottom=8pt, left=10pt, right=10pt,
  },
  tarjetaIcono/.style={
    enhanced, breakable,
    arc=6pt, outer arc=6pt,
    colback=white, colframe=grisLinea, boxrule=0.8pt,
    fonttitle=\bfseries\small, coltitle=azulNoche,
    borderline west={3pt}{0pt}{cyanNeon},
    drop fuzzy shadow=azulNoche!25!white,
    top=6pt, bottom=6pt, left=10pt, right=10pt,
  },
  cajaTitulo/.style={
    enhanced, breakable,
    arc=6pt, outer arc=6pt,
    colback=azulNoche!10!grisPapel, colframe=azulNoche,
    boxrule=1pt,
    fonttitle=\bfseries, coltitle=white,
    attach boxed title to top left={yshift=-3mm, xshift=6mm},
    boxed title style={
      colback=azulNoche, colframe=azulNoche,
      arc=4pt, boxrule=0pt,
    },
    drop fuzzy shadow=azulNoche!25!white,
    top=8pt, bottom=8pt, left=10pt, right=10pt,
  },
  cajaContenido/.style={
    enhanced, breakable,
    arc=5pt, outer arc=5pt,
    colback=grisPapel, colframe=grisLinea,
    boxrule=0.8pt,
    fonttitle=\bfseries\small, coltitle=azulNoche,
    top=6pt, bottom=6pt, left=10pt, right=10pt,
  },
  cajaFortaleza/.style={
    enhanced, breakable,
    arc=6pt, outer arc=6pt,
    colback=verdeSignal!8!white, colframe=verdeSignal,
    boxrule=1pt,
    borderline west={4pt}{0pt}{verdeSignal},
    fonttitle=\bfseries\small, coltitle=verdeSignal!70!black,
    drop fuzzy shadow=verdeSignal!20!white,
    top=6pt, bottom=6pt, left=10pt, right=10pt,
  },
  cajaMejora/.style={
    enhanced, breakable,
    arc=6pt, outer arc=6pt,
    colback=naranjaVivo!8!white, colframe=naranjaVivo,
    boxrule=1pt,
    borderline west={4pt}{0pt}{naranjaVivo},
    fonttitle=\bfseries\small, coltitle=naranjaVivo!80!black,
    drop fuzzy shadow=naranjaVivo!20!white,
    top=6pt, bottom=6pt, left=10pt, right=10pt,
  },
  cajaRecomendacion/.style={
    enhanced, breakable,
    arc=6pt, outer arc=6pt,
    colback=cyanNeon!6!white, colframe=cyanNeon!60!azulNoche,
    boxrule=1pt,
    borderline west={4pt}{0pt}{cyanNeon},
    fonttitle=\bfseries\small, coltitle=azulNoche,
    drop fuzzy shadow=azulNoche!20!white,
    top=6pt, bottom=6pt, left=10pt, right=10pt,
  },
}

% ── Ícono + texto (encabezados de sección y elementos) ────────────────────────
\newcommand{\iconotexto}[2]{\textcolor{cyanNeon}{\faIcon{#1}}~#2}

% ── Listados de código: estilo dark + auto-envoltura ─────────────────────────
\lstdefinestyle{estiloCodigo}{
    backgroundcolor=\color{codigoFondo},
    basicstyle=\ttfamily\small\color{white},
    commentstyle=\color{codigoComentario}\itshape,
    stringstyle=\color{codigoCadena},
    keywordstyle=\color{codigoKeyword}\bfseries,
    numberstyle=\tiny\color{codigoComentario},
    numbers=left,
    numbersep=10pt,
    showstringspaces=false,
    breaklines=true,
    breakatwhitespace=false,
    tabsize=4,
    frame=none,
    captionpos=b,
    aboveskip=6pt,
    belowskip=4pt,
    xleftmargin=14pt,
    framexleftmargin=14pt,
    literate=
      {á}{{\'a}}1 {é}{{\'e}}1 {í}{{\'i}}1 {ó}{{\'o}}1 {ú}{{\'u}}1
      {Á}{{\'A}}1 {É}{{\'E}}1 {Í}{{\'I}}1 {Ó}{{\'O}}1 {Ú}{{\'U}}1
      {ñ}{{\~n}}1 {Ñ}{{\~N}}1 {ü}{{\"u}}1 {Ü}{{\"U}}1
      {¿}{{?`}}1 {¡}{{!`}}1
}
\lstdefinestyle{estiloPython}{
    style=estiloCodigo,
    language=Python,
}
\lstset{style=estiloCodigo}
\BeforeBeginEnvironment{lstlisting}{%
  \begin{tcolorbox}[
    enhanced, breakable,
    arc=6pt, outer arc=6pt,
    colback=codigoFondo, colframe=azulNoche,
    boxrule=1pt,
    drop fuzzy shadow=azulNoche!25!white,
    top=2pt, bottom=2pt, left=0pt, right=0pt,
  ]%
}
\AfterEndEnvironment{lstlisting}{\end{tcolorbox}}

% ── Cajas de contenido temáticas ──────────────────────────────────────────────
\newtcolorbox{cajaRecuerda}{
    enhanced, breakable,
    arc=6pt, outer arc=6pt,
    colback=azulNoche!10!grisPapel,
    colframe=azulNoche, boxrule=1pt,
    borderline west={4pt}{0pt}{cyanNeon},
    fonttitle=\bfseries\small\color{white},
    title={\faIcon{info-circle}~Recuerda},
    attach boxed title to top left={yshift=-3mm, xshift=6mm},
    boxed title style={colback=azulNoche, arc=4pt, boxrule=0pt},
    drop fuzzy shadow=azulNoche!25!white,
    left=10pt, right=10pt, top=8pt, bottom=8pt,
}

\newtcolorbox{cajaConcepto}{
    enhanced, breakable,
    arc=6pt, outer arc=6pt,
    colback=verdeSignal!8!white,
    colframe=verdeSignal, boxrule=1pt,
    borderline west={4pt}{0pt}{verdeSignal},
    fonttitle=\bfseries\small\color{white},
    title={\faIcon{lightbulb}~Concepto clave},
    attach boxed title to top left={yshift=-3mm, xshift=6mm},
    boxed title style={colback=verdeSignal!80!black, arc=4pt, boxrule=0pt},
    drop fuzzy shadow=verdeSignal!20!white,
    left=10pt, right=10pt, top=8pt, bottom=8pt,
}

\newtcolorbox{cajaEjemplo}{
    enhanced, breakable,
    arc=6pt, outer arc=6pt,
    colback=cyanNeon!5!white,
    colframe=cyanNeon!70!azulNoche, boxrule=1pt,
    borderline west={4pt}{0pt}{cyanNeon},
    fonttitle=\bfseries\small\color{fondoOscuro},
    title={\faIcon{code}~Ejemplo},
    attach boxed title to top left={yshift=-3mm, xshift=6mm},
    boxed title style={colback=cyanNeon, arc=4pt, boxrule=0pt},
    drop fuzzy shadow=azulNoche!20!white,
    left=10pt, right=10pt, top=8pt, bottom=8pt,
}

\newtcolorbox{cajaEjercicio}{
    enhanced, breakable,
    arc=6pt, outer arc=6pt,
    colback=azulMedio!7!white,
    colframe=azulMedio, boxrule=1pt,
    borderline west={4pt}{0pt}{azulMedio},
    fonttitle=\bfseries\small\color{white},
    title={\faIcon{pencil-alt}~Ejercicio},
    attach boxed title to top left={yshift=-3mm, xshift=6mm},
    boxed title style={colback=azulMedio, arc=4pt, boxrule=0pt},
    drop fuzzy shadow=azulNoche!20!white,
    left=10pt, right=10pt, top=8pt, bottom=8pt,
}

\newtcolorbox{cajaReto}{
    enhanced, breakable,
    arc=6pt, outer arc=6pt,
    colback=naranjaVivo!6!white,
    colframe=naranjaVivo, boxrule=1pt,
    borderline west={4pt}{0pt}{naranjaVivo},
    fonttitle=\bfseries\small\color{white},
    title={\faIcon{fire}~Reto},
    attach boxed title to top left={yshift=-3mm, xshift=6mm},
    boxed title style={colback=naranjaVivo, arc=4pt, boxrule=0pt},
    drop fuzzy shadow=naranjaVivo!20!white,
    left=10pt, right=10pt, top=8pt, bottom=8pt,
}

% ── Indicadores de dificultad ─────────────────────────────────────────────────
\newcommand{\facil}{\textcolor{verdeSignal}{\faIcon{circle}\,\textbf{Fácil}}}
\newcommand{\intermedio}{\textcolor{naranjaVivo}{\faIcon{circle}\,\textbf{Intermedio}}}
\newcommand{\dificil}{\textcolor{rojoAlerta}{\faIcon{circle}\,\textbf{Difícil}}}

% ─────────────────────────────────────────────────────────────────────────────
\renewcommand{\iconoBanda}{flask}
\fancyhead[L]{\footnotesize\color{grisTexto}\faIcon{flask}~LABORATORIO\_I\_FISICA}
\fancyhead[R]{\footnotesize\color{grisTexto}Guía del estudiante}
\hypersetup{pdftitle={Laboratorio I de Física},pdfauthor={Prof. César Rodríguez}
  pdfsubject={LABORATORIO\_I\_FISICA}}

\begin{document}
\bandaTitulo[fondoOscuro]{Guía del estudiante}{Laboratorio I de Física --- cómo se trabaja y se entrega}

\begin{center}\small\color{grisTexto}Prof. César Rodríguez \textbar{} Universidad de Oriente, Núcleo de Sucre, Departamento de Física \textbar{} 6 de noviembre de 2026\end{center}

\section{\iconotexto{info-circle}{Cómo funciona el curso}}

Las cuatro primeras semanas son talleres de teoría de errores y análisis gráfico: se aprende a medir, a propagar incertidumbres y a ajustar antes de enfrentarse a un montaje. A partir de la semana 5 cada sesión es una práctica experimental completa, con montaje, tabla de datos, análisis y entrega. Cada semana el grupo rota los roles y cierra con un informe que otro equipo podría reproducir.

\section{\iconotexto{users}{Grupos y roles}}

El curso trabaja con 2 grupos. Dos grupos fijos durante el semestre, de tres a cuatro integrantes cada uno. Los roles rotan al inicio de cada fase para que todos los integrantes pasen por la medición, el análisis, la documentación y la verificación.

\begin{itemize}[leftmargin=1.4em,itemsep=2pt,topsep=3pt]
  \item Medición: Prepara el montaje, verifica el cero y el rango de cada instrumento, toma las lecturas y anota la resolución con que fueron leídas. Aporta: La tabla de datos crudos de la semana y el registro de condiciones ambientales e instrumentales.
  \item Análisis: Propaga incertidumbres, linealiza el modelo, ajusta por mínimos cuadrados y extrae los parámetros físicos. Aporta: El script de análisis, las figuras con unidades y el presupuesto de incertidumbres del resultado final.
  \item Documentación: Redacta el informe, integra tablas y figuras, cita las fuentes y mantiene el repositorio ordenado según la estructura acordada. Aporta: El informe de la semana y el README de la carpeta semanal.
  \item Verificación: Contrasta el resultado con el modelo y con una medición independiente, busca errores sistemáticos y firma la trazabilidad de los datos. Aporta: La comparación con el valor de referencia y la lista de desvíos con su causa probable.
\end{itemize}

Los roles giran una posición en sentido horario en cuatro rotaciones (semanas 1-2, 3-4, 5-7 y 8-10), de modo que al cierre de la última cada integrante haya ocupado los cuatro roles.

\textbf{Asignación de roles por rotación (idéntica en cada grupo)}

\begin{tabularx}{\linewidth}{@{}l X X X X@{}}
\toprule
Rotación & Integrante 1 & Integrante 2 & Integrante 3 & Integrante 4 \\
\midrule
Semanas 1-2 & Medición & Análisis & Documentación & Verificación \\
Semanas 3-4 & Análisis & Documentación & Verificación & Medición \\
Semanas 5-7 & Documentación & Verificación & Medición & Análisis \\
Semanas 8-10 & Verificación & Medición & Análisis & Documentación \\
\bottomrule
\end{tabularx}

Cada grupo crea un repositorio privado con el docente como colaborador de lectura. El repositorio es la entrega: se revisa al cierre de cada semana y es la evidencia del trabajo del semestre.

\section{\iconotexto{code-branch}{Portafolio del grupo y entregas}}

La entrega se hace en GitHub. Un repositorio privado por grupo, con el docente como colaborador de lectura. El repositorio es la entrega; no se reciben informes por correo ni impresos, salvo indicación expresa del docente.

El docente revisa el repositorio al finalizar cada semana y devuelve observaciones antes de la semana siguiente.

Corte de cada semana: Domingo a las 23:59 de la semana correspondiente. Lo entregado después del corte se corrige con la penalización prevista en las políticas.

\textbf{Estructura del repositorio del grupo}

\begin{itemize}[leftmargin=1.4em,itemsep=2pt,topsep=3pt]
  \item README.md con integrantes, roles por rotación y tabla de semanas entregadas.
  \item semana\_NN/informe.pdf con el informe de la práctica o el taller resuelto.
  \item semana\_NN/datos/ con los CSV crudos, uno por ensayo, sin editar.
  \item semana\_NN/analisis/ con los scripts que reproducen cada figura.
  \item semana\_NN/figuras/ con las gráficas en PDF o PNG generadas por los scripts.
  \item semana\_NN/README.md con lo que se hizo, lo que no se pudo verificar y por qué.
\end{itemize}

\textbf{Reglas de entrega}

\begin{itemize}[leftmargin=1.4em,itemsep=2pt,topsep=3pt]
  \item Los datos crudos no se editan nunca; toda corrección vive en el script de análisis.
  \item Cada cifra del informe debe salir de una tabla de datos o de un script entregado.
  \item Los scripts se ejecutan de punta a punta sin pasos manuales intermedios.
  \item Ninguna tabla se entrega sin unidades ni sin la resolución del instrumento.
  \item El informe es grupal, pero las conclusiones se firman individualmente.
  \item Un resultado sin incertidumbre no se acepta como resultado.
\end{itemize}

\begin{tcolorbox}[cajaMejora]\textbf{Antes de empezar:} trabaja solo con el material asignado a tu mesa y con los instrumentos cuyo uso haya autorizado el docente. Si un montaje no coincide con la figura o un instrumento no responde, detente y consulta; un resultado dudoso se reporta, no se corrige a mano.
\end{tcolorbox}

\section{\iconotexto{cube}{Puesta en marcha del laboratorio}}

Laboratorio físico de mesa. Cada grupo dispone de un juego de instrumentos y de una estación de análisis con Python para procesar sus propios datos.

Mesas con regla milimetrada, calibrador, tornillo micrométrico, balanza, probeta, cronómetro, termómetro digital, soportes, resortes, masas calibradas y cuerpos sólidos de distintos materiales y tamaños.

\begin{itemize}[leftmargin=1.4em,itemsep=2pt,topsep=3pt]
  \item No se energiza ningún equipo sin revisión previa del docente.
  \item Los líquidos se manipulan lejos de las balanzas y de las conexiones eléctricas.
  \item Cada instrumento se devuelve limpio, seco y en su estuche al terminar la sesión.
  \item El material roto se reporta de inmediato; ocultarlo agrava la sanción.
  \item Ningún dato se inventa ni se completa de memoria después de la sesión.
\end{itemize}

\section{\iconotexto{code}{Contrato de datos de los scripts}}

\begin{cajaConcepto}Todo script entregado imprime \textbf{un solo objeto JSON} por la salida estándar y devuelve un diccionario con estas claves. \texttt{status} vale \texttt{ok} o \texttt{error}; si es \texttt{error}, \texttt{errores} explica el motivo y el programa sale con código distinto de cero.
\begin{lstlisting}[language=Python,style=estiloPython]
{
  "status": "ok",
  "documento": "semana_05",
  "magnitud": "volumen",
  "valor": 12.43,
  "unidad": "cm^3",
  "incertidumbre": 0.08,
  "datos": "datos/volumen.csv",
  "figuras": ["figuras/volumen_metodos.pdf"],
  "modelo": "V = pi d^2 h / 4",
  "parametros": {"d": 2.00, "h": 3.96}
}

\end{lstlisting}
\end{cajaConcepto}

\begin{itemize}[leftmargin=1.4em,itemsep=2pt,topsep=3pt]
  \item El script parte de los CSV crudos y termina en las figuras del informe, sin pasos manuales.
  \item Toda cifra del informe se calcula en el script; nada se teclea a mano.
  \item Las unidades se declaran en el CSV y en los ejes de cada figura.
  \item El ajuste reporta parámetro, incertidumbre, coeficiente de determinación y residuales.
  \item Sin credenciales ni rutas absolutas: el script corre desde la carpeta de la semana.
\end{itemize}

\section{\iconotexto{clipboard-check}{Evaluación}}

\begin{tabularx}{\linewidth}{@{}X r@{}}
\toprule
Concepto & Peso (\%)\\
\midrule
Talleres y prácticas (semanas 1 a 10) & 60,0 \\
Controles de fase (semanas 4 y 7) & 20,0 \\
Evaluación final práctica y portafolio (semana 10) & 20,0 \\
\bottomrule
\end{tabularx}

Cada informe se corrige con la rúbrica del curso, que se entrega junto con el enunciado. La nota es la suma de los criterios, no un promedio: un criterio sin evidencia vale cero aunque el resultado numérico sea correcto.

\section{\iconotexto{balance-scale}{Seguridad, honestidad académica y uso del laboratorio}}

\begin{itemize}[leftmargin=1.4em,itemsep=2pt,topsep=3pt]
  \item Seguridad: no se energiza ningún equipo ni se manipula vidrio caliente sin autorización y sin el equipo de protección indicado.
  \item Honestidad académica: los datos se toman en la sesión y se declaran tal como se obtuvieron; fabricar o alterar datos anula la entrega y se reporta.
  \item Puntualidad: la entrega cierra el domingo a las 23:59; lo entregado después se corrige con la penalización prevista.
  \item Cuidado del material: cada instrumento se devuelve limpio y en su estuche; el material roto se reporta de inmediato.
  \item Trabajo en equipo: todos los integrantes firman las conclusiones y deben poder explicar cualquier parte del informe.
\end{itemize}

\section{\iconotexto{calendar-alt}{Calendario resumido}}

\begin{tabularx}{\linewidth}{@{}r p{4.1cm} X@{}}
\toprule
Sem. & Tema & Entregable\\
\midrule
1 & La medición, la apreciación y el error & Taller resuelto con la tabla de diez lecturas, el cálculo completo de la media y las incertidumbres, la expresión final del resultado y un párrafo que justifique el instrumento elegido. \\
2 & Propagación de incertidumbres y notación científica & Taller resuelto con las mediciones, la propagación paso a paso, ambos volúmenes expresados en notación científica y la conclusión sobre si los métodos coinciden dentro de la incertidumbre. \\
3 & Graficación, linealización y lectura de una recta & Taller resuelto con las dos gráficas, la recta trazada, el cálculo de la pendiente con unidades, el valor de $g$ y la interpretación del intercepto. \\
4 & Mínimos cuadrados y estructura del informe & Taller integrador con la tabla completa, el ajuste por mínimos cuadrados, la gráfica con residuales, la constante elástica y un informe corto de dos páginas con la estructura del curso. \\
5 & Volumen de un sólido por dos métodos & Informe con la tabla de datos crudos, ambos volúmenes con su incertidumbre, la comparación entre métodos y la explicación de la discrepancia. \\
6 & Densidad de un cuerpo sólido & Informe con la tabla de datos, las densidades por ambos métodos con su incertidumbre, la identificación propuesta del material y el análisis de las fuentes de error. \\
7 & Centro de masa de un sistema de dos cuerpos & Informe con la tabla de los tres montajes, los centros de masa calculados con y sin varilla, la comparación con el medido y la conclusión sobre la validez de despreciar la masa de la varilla. \\
8 & Ley de Hooke y constante elástica & Informe con la tabla de datos, la gráfica con el ajuste y los residuales, la constante elástica con su incertidumbre y el rango de validez de la ley de Hooke. \\
9 & Péndulo simple y determinación de la gravedad & Informe con la tabla completa, la gráfica linealizada, el ajuste, el valor de $g$ con su incertidumbre, la corrección de longitud y el análisis de los datos anómalos. \\
10 & Enfriamiento de un cuerpo sólido & Informe final con las tablas de ambos cuerpos, las gráficas linealizadas, las constantes de enfriamiento con su incertidumbre, la comparación de los tiempos característicos y la discusión de las limitaciones del modelo. \\
\bottomrule
\end{tabularx}

\section{\iconotexto{paperclip}{Prefijos del Sistema Internacional}}

Los prefijos se aplican a la unidad y no cambian el número de cifras significativas. Al convertir, es preferible pasar primero a la unidad base y después aplicar el prefijo, para no perder la cuenta de los decimales.

\begin{center}
\begin{minipage}{\linewidth}
\centering
{\small
\textbf{Prefijos de uso frecuente en el laboratorio}\par\smallskip
\begin{tabular}{@{}ccc@{}}
\toprule
Prefijo & Símbolo & Factor \\
\midrule
kilo & k & $10^{3}$ \\
mili & m & $10^{-3}$ \\
micro & $\mu$ & $10^{-6}$ \\
nano & n & $10^{-9}$ \\
mega & M & $10^{6}$ \\
centi & c & $10^{-2}$ \\
\bottomrule
\end{tabular}
}
\end{minipage}
\end{center}

\section{\iconotexto{paperclip}{Densidades de materiales comunes}}

Valores de referencia a temperatura ambiente. La identificación de un material por su densidad es ambigua cuando dos materiales tienen valores cercanos; en ese caso se necesita otra propiedad, como el color o la respuesta a un imán.

\begin{center}
\begin{minipage}{\linewidth}
\centering
{\small
\textbf{Densidad de sólidos y líquidos de uso frecuente}\par\smallskip
\begin{tabular}{@{}cc@{}}
\toprule
Material & $\rho$ (g/cm$^3$) \\
\midrule
Aluminio & 2,70 \\
Vidrio & 2,50 \\
Hierro & 7,87 \\
Cobre & 8,96 \\
Plomo & 11,34 \\
Agua a 20 C & 0,998 \\
Glicerina & 1,26 \\
\bottomrule
\end{tabular}
}
\end{minipage}
\end{center}

\section{\iconotexto{paperclip}{Constantes y valores de referencia}}

Se usan para comparar los resultados experimentales, no para sustituir la medición. Toda comparación debe expresarse como una diferencia porcentual y discutirse frente a la incertidumbre del resultado.

\begin{center}
\begin{minipage}{\linewidth}
\centering
{\small
\textbf{Valores de referencia del curso}\par\smallskip
\begin{tabular}{@{}ccc@{}}
\toprule
Magnitud & Valor & Unidad \\
\midrule
Aceleración de la gravedad & 9,80 & m/s$^2$ \\
Densidad del agua a 4 C & 1,000 & g/cm$^3$ \\
Densidad del agua a 20 C & 0,998 & g/cm$^3$ \\
Calor específico del aluminio & 0,900 & J/(g K) \\
\bottomrule
\end{tabular}
}
\end{minipage}
\end{center}
\end{document}
